\section{R2: critical-window Gaussian approximation}
\label{sec:R2}

\Cref{thm:R1} leaves the intermediate window
$\snet\in(\critrate/\sqrt2,\sqrt2\,\critrate)$ uncovered. R2 fills it with
a Gaussian-CDF approximation of the true-decoder success probability,
obtained by applying a martingale Berry--Esseen theorem
(\cite{hallheyde1980}, Theorem~3.6) to the centred signal martingale of
\Cref{lem:signal_floor} on a high-probability truncation event from
\Cref{lem:orthogonality}.

\begin{theorem}[Gaussian-CDF approximation in the critical window]
\label{thm:R2}
Suppose \Cref{ass:descent,ass:bounded_architecture} hold, with the mild
incoherence condition \Cref{eq:R1_incoherence} ($\incoh_0<\rho_0/R_U$). Fix
$d\ge1$, $|\Vocab|\ge3$, and a horizon $T_{\max}\ge1$. For every $\snet$ in
the critical window $(\critrate/\sqrt2,\sqrt2\,\critrate)$, the true-decoder
success probability satisfies
\begin{align}\label{eq:R2_phi}
   \Pr\bigl[\, \Dec(x_{T_{\max}})=\corr \,\bigr]
   \;=\;
   \Phi(z) \;+\; O\!\bigl(\beerr\bigr),
\end{align}
where $\Phi$ is the standard normal CDF, the standardised threshold is
\begin{align}\label{eq:R2_z}
   z
   \;\coloneqq\;
   \frac{(\snet-\critrate)\,\driftc\,\rho_0\,(\rho_0-\incoh_0 R_U)}{\sigma_{T_{\max}}},
   \qquad
   \sigma_{T_{\max}}^2 \;\le\; \frac{4\,R_U^2 M^2 e^{2S}}{T_{\max}\,d},
\end{align}
is the conditional variance of the centred margin martingale at horizon
$T_{\max}$ (\Cref{lem:signal_floor}, \Cref{lem:quad_variation}), the drift
coefficient $\driftc\rho_0(\rho_0-\incoh_0 R_U)$ is the margin-drift slope
of \Cref{thm:R1} (under the mild incoherence \Cref{eq:R1_incoherence}), and
the Berry--Esseen rate is
\begin{align}\label{eq:R2_eps}
   \beerr
   \;=\;
   O\!\left(
      \frac{e^{S}}{\sqrt{T_{\max}}}
      \;+\;
      \sqrt{\frac{\log|\Vocab|}{T_{\max}\,d}}
   \right).
\end{align}
In particular $z=0$ at $\snet=\critrate$, so $\Phi(0)=\tfrac12$: at the
exact threshold the success probability is $\tfrac12+O(\beerr)$. Since
$z\propto(\snet-\critrate)\sqrt{T_{\max}\,d}$, as $T_{\max}d\to\infty$ the
curve sharpens to the step function of \Cref{thm:R1}.
\end{theorem}

\begin{proof}
The proof applies the Hall--Heyde martingale Berry--Esseen theorem
(\Cref{thm:hallheyde_be}) to the \emph{correct-token logit}, a martingale in
the \emph{fixed} direction $W_U^{\corr}$, and treats the incorrect side as a
deterministic-on-the-good-event threshold pinned between
\Cref{lem:incorrect_max} (upper) and \Cref{lem:incorrect_max_lower} (lower).
As in \Cref{thm:R1} the comparison is on $\tilde x_{T_{\max}}$, the rescaling
to $x_{T_{\max}}$ preserving $\operatorname{sign}\Margin$.

\smallskip\noindent\textbf{Union-bound paragraph.} The $O(\beerr)$ error
absorbs an explicit union bound over three events: (1) the truncation event
$\Ecal_{\mathrm{trunc}}$ of \Cref{lem:orthogonality} applied to each
\emph{correct-direction} per-step projection $\inner{W_U^{\corr}}{\xi_t}/R_U$,
failure $\le O(T_{\max}e^{-cd})$ via a union over the $T_{\max}$ steps; (2)
the incorrect-max event $\Ecal_{\mathrm{inc}}$ of
\Cref{lem:incorrect_max,lem:incorrect_max_lower} pinning the threshold
$\Theta_{T_{\max}}\coloneqq\max_{a\neq\corr}\inner{W_U^a}{\tilde x_{T_{\max}}}$,
failure $\le(|\Vocab|-1)^{-1}$; (3) the Hall--Heyde output, deterministic on
$\Ecal_{\mathrm{trunc}}$, no extra budget. The total
$O(T_{\max}e^{-cd})+(|\Vocab|-1)^{-1}$ is dominated by the polynomial rate
\Cref{eq:R2_eps}.

\smallskip\noindent\textbf{Step 1 (decompose into a fixed-direction
martingale plus a pinned threshold).} Write the correct logit as its
$\Fcal_0$-conditional drift plus a centred martingale in the \emph{fixed}
direction $W_U^{\corr}$,
\begin{align*}
   \inner{W_U^{\corr}}{\tilde x_{T_{\max}}}
   \;=\; \underbrace{\norm{W_U^{\corr}}_2\, D}_{\text{drift }\mu_{T_{\max}}}
       \;+\; \underbrace{\textstyle\sum_{t} D_t}_{\text{martingale}},
   \qquad
   D_t=w_{T,t}\norm{W_U^{\corr}}_2\inner{\rowhat}{\xi_t},
\end{align*}
where $D=\inner{\rowhat}{\E\tilde x_{T_{\max}}}=\snet\driftc\rho_0\cdot(1+o(1))$
at leading order (\Cref{lem:signal_floor} Step~2; the drift is exactly
$\snet\,\norm{W_U^{\corr}}_2\,\E_0[\sum_k w_{T,k}\E[\abs{\proj_k}\mid\Fcal_{k-1}]]$
by the net-drift identity, a fixed-direction quantity with no
data-dependence). On $\Ecal_{\mathrm{inc}}$ the incorrect threshold is pinned,
$\Theta_{T_{\max}}\in[\,\sigma_{T_{\max}}\sqrt{2\kappa(1-\incoh_0)\log(|\Vocab|-1)}-O(\snet),\ \incoh_0 R_U\abs D+\sigma_{T_{\max}}\sqrt{2\log(|\Vocab|-1)}\,]$,
both endpoints equal to $\critrate\,\norm{W_U^{\corr}}_2\driftc\rho_0(\rho_0-\incoh_0 R_U)/\rho_0$
up to the $\sqrt2$ slack by the definition \Cref{eq:critical_rate} of
$\critrate$ (this pinning is where the $\sqrt2$ slack and the $(1-\incoh_0)$
factor of \Cref{thm:R1} are spent). Hence, deterministically on
$\Ecal_{\mathrm{inc}}$,
$\{\Margin>0\}=\{\inner{W_U^{\corr}}{\tilde x_{T_{\max}}}>\Theta_{T_{\max}}\}=\{\sum_t D_t/\sigma_{T_{\max}}>-z+O(\beerr)\}$
with $z$ of \Cref{eq:R2_z} and the $O(\beerr)$ absorbing the threshold width.

\smallskip\noindent\textbf{Step 2 (Berry--Esseen for the fixed-direction
martingale).} The differences $(D_t)$ form a square-integrable martingale
array in the fixed direction $\rowhat$ (\Cref{ass:bounded_architecture}(2)),
with terminal functional independent of the data. Its conditional variance
is, by the isotropy bound \Cref{eq:noise_isotropy} and
$\norm{W_U^{\corr}}_2\le R_U$,
$V_{T_{\max}}=\sum_t\E[D_t^2\mid\Fcal_{t-1}]\le S_{T_{\max}}R_U^2\sigma_{\max}^2/d\le4R_U^2M^2e^{2S}/(T_{\max}d)=\sigma_{T_{\max}}^2$
(\Cref{lem:quad_variation}, $\sigma_{\max}\le2M$). On $\Ecal_{\mathrm{trunc}}$
the per-step projection is sub-Gaussian with truncated scale $O(1/\sqrt d)$,
so the conditional third moment is
$\rho_{T_{\max}}=\sum_t\E[|D_t|^3\mid\Fcal_{t-1}]=O\bigl(\sum_t w_{T,t}^3 R_U^3 M^3 d^{-3/2}\bigr)=O(R_U^3 M^3 e^{3S}d^{-3/2}T_{\max}^{-2})$,
the $d^{-3/2}$ coming from the truncated tail (sharper than the
deterministic per-step bound). Hence the Berry--Esseen rate of
\Cref{thm:hallheyde_be} is
$\rho_{T_{\max}}/V_{T_{\max}}^{3/2}=O(e^S/\sqrt{T_{\max}})$. Applying
\Cref{thm:hallheyde_be} to $(\sum_t D_t)/\sigma_{T_{\max}}$,
\begin{align*}
   \Bigl|\Pr\bigl[\textstyle\sum_t D_t\le x\,\sigma_{T_{\max}}\bigr]-\Phi(x)\Bigr|
   \;\le\; C\,\frac{\rho_{T_{\max}}}{V_{T_{\max}}^{3/2}}
   \;=\; O\!\bigl(e^S/\sqrt{T_{\max}}\bigr).
\end{align*}

\smallskip\noindent\textbf{Step 3 (invert).} With $x=-z+O(\beerr)$,
\begin{align*}
   \Pr[\Margin>0]
   \;=\; \Pr\bigl[\textstyle\sum_t D_t/\sigma_{T_{\max}}>-z+O(\beerr)\bigr]
   \;=\; 1-\Phi(-z)+O(\beerr)
   \;=\; \Phi(z)+O(\beerr),
\end{align*}
by symmetry of $\Phi$. The threshold-pinning term
$\sqrt{\log|\Vocab|/(T_{\max}d)}$ from the incorrect-max side
(\Cref{lem:incorrect_max}) adds to the Hall--Heyde $e^S/\sqrt{T_{\max}}$ to
give the rate \Cref{eq:R2_eps}. This completes the proof.
\end{proof}

\begin{corollary}[Boundary value at the threshold]
\label{cor:R2_boundary}
At $\snet=\critrate$, $z=0$ and
$\Pr[\Dec(x_{T_{\max}})=\corr]=\tfrac12+O(\beerr)$.
\end{corollary}

\begin{proof}
Substituting $\snet=\critrate$ into \Cref{eq:R2_z} makes the numerator
$(\snet-\critrate)=0$, so $z=0$; then $\Phi(0)=\tfrac12$ in
\Cref{eq:R2_phi} gives $\Pr[\Dec(x_{T_{\max}})=\corr]=\tfrac12+O(\beerr)$.
This completes the proof.
\end{proof}

\begin{corollary}[Transition width, absolute scale]
\label{cor:R2_width}
Under the hypotheses of \Cref{thm:R2}, the success probability rises from
$0.1$ to $0.9$ as $\snet$ sweeps an \emph{absolute} window in $\snet$ of
width
\begin{align}\label{eq:R2_width_abs}
   w_{\mathrm{abs}} \;=\; \Theta\!\left(\frac{1}{\sqrt{T_{\max}\,d}}\right),
\end{align}
centred at $\critrate$.
\end{corollary}

\begin{proof}
By \Cref{eq:R2_phi} the success probability is $\Phi(z)+O(\beerr)$, which
rises from $0.1$ to $0.9$ as $z$ sweeps
$[\Phi^{-1}(0.1),\Phi^{-1}(0.9)]\approx[-1.28,1.28]$. Inverting the linear
map \Cref{eq:R2_z} between $z$ and $\snet$, the corresponding $\snet$-width
is $w_{\mathrm{abs}}=2\cdot1.28\,\sigma_{T_{\max}}/[\driftc\rho_0(\rho_0-\incoh_0 R_U)]$;
substituting $\sigma_{T_{\max}}\le2R_U M e^{S}/\sqrt{T_{\max}d}$ gives
$w_{\mathrm{abs}}=\Theta(1/\sqrt{T_{\max}d})$, which is
\Cref{eq:R2_width_abs}. This completes the proof.
\end{proof}

\begin{remark}[Absolute vs.\ relative width: do not conflate]
\label{rem:R2_width_relative}
\Cref{cor:R2_width} reports the \emph{absolute} width in $\snet$, the
headline scale $\Theta(1/\sqrt{T_{\max}d})$. A second, distinct, notion is
the width \emph{relative} to the critical rate $\critrate$ itself:
$w_{\mathrm{rel}}=w_{\mathrm{abs}}/\critrate=\Theta\bigl((\log|\Vocab|)^{-1/2}\bigr)$,
since the $\sqrt{1/(T_{\max}d)}$ factors cancel in the ratio. These two
widths are \emph{not} the same quantity and must not be equated: the
absolute width shrinks as $T_{\max}d\to\infty$ (the transition sharpens),
whereas the relative width depends only on $\log|\Vocab|$ and is
$\Theta(1)$ in $d$ and $T_{\max}$. The headline scaling statement of the
paper is the absolute width \Cref{eq:R2_width_abs}; the relative width is
recorded here as a separate empirical signature.
\end{remark}
